\documentclass[10pt]{beamer}

\usetheme{Madrid}
\usecolortheme{default}
\setbeamertemplate{navigation symbols}{}

\usepackage{amsmath,amssymb}
\usepackage{mathtools}

\title[Power-Spectrum Constraints]{Two Implementations of Power-Spectrum Constraints}
\author{}
\date{}

\begin{document}

\begin{frame}
  \titlepage
\end{frame}

\begin{frame}{Where the constraint enters}
  \begin{itemize}
    \item In \texttt{compsep\_planck.py}, \texttt{ST\_COMPUTE\_PS=1} sets
    \[
      \texttt{compute\_PS=True}
    \]
    when building the \texttt{ST\_Operator}.

    \item The backend then appends a separate tensor \texttt{PS} to the flattened
    statistics:
    \[
      \texttt{mean},\ \texttt{var},\ S_1,\ldots,S_4,\ \texttt{PS}.
    \]

    \item The target and running maps are compared by the same squared
    Euclidean loss used for the other statistics.
  \end{itemize}
\end{frame}

\begin{frame}{Planck backend: channel pairs}
  In the 5-channel phase,
  \[
    X = (d_Q,\ n_Q,\ d_U,\ n_U,\ d_I).
  \]
  The power-spectrum operator uses the same upper-triangular pair mask as the
  scattering statistics:
  \[
  \begin{pmatrix}
  1&0&1&0&1\\
  0&1&0&0&0\\
  0&0&1&0&1\\
  0&0&0&1&0\\
  0&0&0&0&1
  \end{pmatrix}.
  \]

  \begin{itemize}
    \item Auto-spectra are computed for every channel.
    \item Cross-spectra are computed for
    \[
      (d_Q,d_U),\quad (d_Q,d_I),\quad (d_U,d_I).
    \]
    \item Signal-noise and noise-auxiliary cross-spectra are not included.
  \end{itemize}
\end{frame}

\begin{frame}{Planck backend: what is computed}
  For each batch element, selected channel pair, and radial bin,
  \[
    P_{c_1c_2}(b)
    =
    \frac{
      \sum_k M_b(k)\,\widehat X_{c_1}(k)
      \overline{\widehat X_{c_2}(k)}
    }{
      \sum_k M_b(k)
    }.
  \]

  \begin{itemize}
    \item \(M_b(k)\) is a smooth radial Fourier-bin mask.
    \item Output shape is
    \[
      N_{\rm batch}\times N_c\times N_c\times N_{\rm bins}.
    \]
    \item Only the selected upper-triangular entries are filled; the redundant
    conjugate entries are left unused.
    \item The result is normalized by reference auto-spectra:
    \[
      P_{c_1c_2}(b) \leftarrow
      \frac{P_{c_1c_2}(b)}
      {\sqrt{P^{\rm ref}_{c_1c_1}(b)}\sqrt{P^{\rm ref}_{c_2c_2}(b)}}.
    \]
  \end{itemize}
\end{frame}

\begin{frame}{Planck backend: non-periodic maps}
  The note \texttt{pbc\_power\_spectrum\_note.pdf} describes the two backend
  branches.

  \begin{itemize}
    \item If \texttt{PBC=True}, the estimate stays in Fourier space.

    \item If \texttt{PBC=False}, the backend first inverse-transforms the
    binned map,
    \[
      X^{(b)}_{c_1}(x)
      =
      \mathcal{F}^{-1}\!\left[M_b(k)\widehat X_{c_1}(k)\right](x),
    \]
    then averages
    \[
      X^{(b)}_{c_1}(x)\overline{X_{c_2}(x)}
    \]
    with a bin-dependent real-space taper/crop mask.

    \item So the intended object is still a radial auto/cross power spectrum;
    the non-periodic branch changes the estimator to reduce boundary artifacts.
  \end{itemize}
\end{frame}

\begin{frame}{\texttt{Scattering2d.py}: what is called power}
  In \texttt{Scattering2d.py}, the main single-field power term is
  \[
    P_{00}(j,\ell)
    =
    \left\langle
      |I\ast\psi_{j,\ell}|^2
    \right\rangle_x .
  \]

  \begin{itemize}
    \item This is wavelet-band energy, not a separately binned radial FFT
    spectrum.
    \item It is included directly in \texttt{for\_synthesis} as
    \[
      \log P_{00}(j,\ell).
    \]
    \item The isotropic vector uses
    \[
      P_{00}^{\rm iso}(j)=\frac{1}{L}\sum_\ell P_{00}(j,\ell).
    \]
    \item A newer extra term \texttt{P00g} also computes energy after filtering
    by Gaussian concentric rings, and appends \(\log P_{00g}\) to the
    non-isotropic synthesis vector.
  \end{itemize}
\end{frame}

\begin{frame}{\texttt{Scattering2d.py}: cross-field terms}
  The two-field routine computes auto powers for each field:
  \[
    P_{00,a}(j,\ell),\qquad P_{00,b}(j,\ell).
  \]

  It also computes the normalized cross-band correlation
  \[
    {\rm Corr}_{00}(j,\ell)
    =
    \frac{
      \left\langle
        \widehat a(k)\overline{\widehat b(k)}
        |\psi_{j,\ell}(k)|^2
      \right\rangle_k
    }{
      \sqrt{P_{00,a}(j,\ell)P_{00,b}(j,\ell)}
    }.
  \]

  \begin{itemize}
    \item Therefore auto and cross information can both appear.
    \item But this interface is pairwise: it takes two fields \(a,b\), not an
    arbitrary \(N_c\times N_c\) channel-pair mask.
    \item The flattened vector contains \(P_{00,a}\), \(P_{00,b}\), and real and
    imaginary parts of \({\rm Corr}_{00}\).
  \end{itemize}
\end{frame}

\begin{frame}{Main difference}
  \begin{columns}[T,onlytextwidth]
    \begin{column}{0.49\textwidth}
      \textbf{Current backend}
      \begin{itemize}
        \item Separate \texttt{PS} block.
        \item Radial Fourier bins \(M_b(k)\).
        \item Multi-channel auto/cross tensor.
        \item Pair selection controlled by one channel-pair mask.
        \item Non-periodic mode uses a real-space taper/crop estimator.
      \end{itemize}
    \end{column}
    \begin{column}{0.49\textwidth}
      \textbf{\texttt{Scattering2d.py}}
      \begin{itemize}
        \item Power terms are embedded in scattering covariance vectors.
        \item Main term is wavelet-band energy \(P_{00}\).
        \item Optional \texttt{P00g} uses Gaussian rings.
        \item Cross-field power appears as normalized pairwise correlation
        \({\rm Corr}_{00}\).
        \item No general multi-channel PS tensor.
      \end{itemize}
    \end{column}
  \end{columns}
\end{frame}

\begin{frame}{Practical takeaway}
  \begin{itemize}
    \item The backend \texttt{st\_compute\_ps} constraint matches radial
    auto/cross Fourier spectra directly.

    \item The \texttt{Scattering2d.py} implementation constrains band energies
    and band correlations attached to the wavelet scattering representation.

    \item Both include auto-power information.

    \item Cross-power differs most:
    \begin{itemize}
      \item backend: selected entries of a full channel-pair cross-spectrum
      tensor;
      \item \texttt{Scattering2d.py}: pairwise normalized cross-correlations
      for two input fields.
    \end{itemize}

    \item The backend is closer to an explicit power-spectrum loss; the older
    code is closer to adding power-like scattering covariance coordinates.
  \end{itemize}
\end{frame}

\end{document}
